% Section 1.2, from lectures/L01/appendix/b_toolchain.md.

\section{The Toolchain}
\label{c1:sec:toolchain}

\subsection{What assembling actually does}
\label{c1:sec:assembling}
An assembler is the least clever program in your toolchain. It reads one line at a time, turns
each instruction into the one or two words that encode it (\secref{c1:sec:encoding}), resolves the
labels, and writes the result out. It does not reorder anything, it does not allocate registers,
and it does not check that what you wrote makes sense. If you write \code{ldi r16, 0x2A}, you get
\code{0xE20A}, and if you write a loop with no exit you get a loop with no exit.

That is the whole appeal. Everything the machine does is something you wrote.

\bookfigure{toolchain}{The toolchain: \code{led.asm} through \code{avra} into a listing, a hex and
a map. The hex and the map together feed the test suite; the hex alone feeds
\code{avrdude}.}{fig:toolchain}

\subsection{\code{avra} as an assembler}
\label{c1:sec:avra}
This book assembles with \code{avra}, which takes AVRASM2, the syntax Microchip Studio assembles,
and never uses the C compiler except in Chapter~6:

\begin{shell}
avra -I /usr/share/avra -I drivers/include -D F_CPU=16000000 \
     -l drivers/build/drivers.lst -m drivers/build/drivers.map \
     -o drivers/build/drivers.hex drivers/build/drivers_unit.asm
\end{shell}

Each of those is load-bearing:

\begin{booktable}{@{}p{10.4em}L@{}}
  \thead{Flag} & \thead{Why} \\
  \midrule
  \mbox{\code{-I /usr/share/avra}} & Where \code{m328Pdef.inc} lives, which is what gives you
    \code{PORTB} and \code{RAMEND}. \\
  \mbox{\code{-I drivers/include}} & Where the course's own \code{.inc} contracts live. \\
  \mbox{\code{-D F_CPU=16000000}} & The clock, as a symbol your own source may use. Nothing in the book
    requires it; Chapter~5's arithmetic is easier to write against a name than a literal. \\
  \code{-l ....lst} & The listing: your source with the encoding beside every line.
    \secref{c1:sec:disasm}. \\
  \code{-m ....map} & The map: every label and its address. \textbf{The tests read this.} \\
  \code{-o ....hex} & Intel hex, which is what the simulator loads. \\
\end{booktable}

You will rarely type that, because \code{make build} does it. Two of the values are overridable
from the environment for the same reason: \code{F_CPU=8000000 make build} assembles against a
different clock, and \code{AVRA_INCLUDE=... make build} is what a non-apt install of \code{avra}
needs, because its device files will not be in \file{/usr/share/avra}.

\textbf{The device is chosen by the file you include, not by a flag.} There is no \code{-mmcu}
here; \code{.include "m328Pdef.inc"} is the whole of that decision, and it is why the line belongs
at the top of every source file you write.

\textbf{\code{avra} has no linker.} It assembles exactly one file, so \file{ci/build.sh} writes a
top-level unit that \code{.include}s every source in \file{drivers/source} and assembles that
instead. You never edit it. When you want to know what was assembled and in what order, read
\file{drivers/build/drivers_unit.asm}; it is one line per source plus two, and generated fresh on
every build.

Two conventions the files in this book follow:
\begin{itemize}
  \item \textbf{\code{;} starts a comment}, to the end of the line.
  \item \textbf{Every label is global already.} \code{avra} exports all of them, so there is no
    \code{.global} to write and no way to forget it. A test that cannot find your subroutine is
    telling you about a spelling or a file that did not assemble, never about a missing
    directive.
\end{itemize}

\subsection{Reading a disassembly}
\label{c1:sec:disasm}
The listing file is the assembler telling you what it produced, line by line, against the source
you gave it. \code{make build} writes one every time:

\begin{shell}
less drivers/build/drivers.lst
\end{shell}

\begin{console}
C:000000 01fc          movw r30, r24            ; Z = the struct
C:000001 e203          ldi  r16, low(PINB + 0x20)
C:000002 8300          std  Z + LED_PIN_REG, r16
\end{console}

The first column is the \textbf{word} address, the second is the encoding, and the rest is your
own line, comment included. That is the source you wrote, the bits it became, and the address
each landed at, all on one line.

To disassemble instead, starting from the machine code and recovering the instructions, which is
the skill this book is actually after, point \code{avr-objdump} at the hex:

\begin{shell}
avr-objdump -D -m avr:5 -b ihex drivers/build/drivers.hex
\end{shell}

\code{-b ihex} is needed because a hex file says nothing about what architecture it is for, and
\code{-m avr:5} supplies the answer. \textbf{The symbol names are gone}, and the output says
\code{.sec1} where an ELF would have said \code{<led_init>}; a hex file carries addresses and
bytes and nothing else, which is exactly why the tests read the map file alongside it.

\begin{console}
Disassembly of section .sec1:

00000000 <.sec1>:
   0:   21 e0           ldi     r18, 0x01       ; 1
   2:   30 e0           ldi     r19, 0x00       ; 0
   4:   38 17           cp      r19, r24
   6:   19 f0           breq    .+6             ; 0xe
   8:   22 0f           add     r18, r18
   a:   33 95           inc     r19
   c:   fb cf           rjmp    .-10            ; 0x4
   e:   82 2f           mov     r24, r18
  10:   08 95           ret
\end{console}

Four columns: the byte address, the encoded words, the instruction, and a comment. Four things in
that listing are worth noticing on your first read.

\textbf{There is one block and it is called \code{.sec1}.} Not \code{shift_bits}, not
\code{shift_bits_loop}: a hex file has no symbol table, so \code{objdump} invents one section name
for the whole run of bytes and labels every branch target with a bare address. The names are in
\file{drivers.map}, which is the other file \code{avra} wrote and the one the tests read. Compare
this against the listing above, where your own label sits on the line that produced it. The
listing knows the names because it is the assembler talking, and the disassembly does not because
it is only ever given bytes.

\textbf{The bytes are little-endian.} \code{ldi r18, 0x01} shows as \code{21 e0}, and the word is
\code{0xE021}. Hold that against \secref{c1:sec:encoding}: opcode \code{1110}, \code{K[7:4]} =
\code{0000}, \code{d} = \code{0010} (so \code{r18}), \code{K[3:0]} = \code{0001}. It checks out,
and reading it in the wrong byte order is a small rite of passage.

\textbf{Addresses here are bytes, not words.} The \code{rjmp} at the bottom goes back to
\code{0x04}, which is word \code{0x02} and so the \emph{third} instruction: words \code{0x00} and
\code{0x01} hold the two \code{ldi}s above it. The vector table is quoted in words and the
disassembly in bytes, and keeping the two apart is a recurring nuisance that this section cannot
make go away.

\textbf{\code{lsl r18} came back as \code{add r18, r18}.} They are the same instruction: shifting
left by one and adding a number to itself do the same thing to the bits and set the same flags,
so the assembler encodes \code{lsl} as \code{add} and the disassembler has no way to know which
you wrote. Several AVR mnemonics are aliases like this, and a disassembly that does not match your
source line for line is usually one of them rather than a mistake.

\subsection{Two images, for two jobs}
\label{c1:sec:images}
\code{make build} produces two, and confusing them wastes an afternoon.

\textbf{\file{drivers/build/drivers.hex}} is the subroutines alone. It has no reset vector and no
vector table; nothing ever runs it from the start. The unit tests set the program counter
straight to a symbol and call one subroutine at a time. This is the image nearly everything in
this book loads, and it is the one \code{make measure} uses by default.

\textbf{\file{drivers/build/app.hex}} is \file{drivers/app/main.asm} assembled together with
those same subroutines: a whole program with a vector table and a main loop. This chapter's
program test and Chapter~3's integration test load this one and let it run, because a program
cannot be called and an interrupt has to arrive. Reach it with \code{make measure IMAGE=app}.

Each comes with a \code{.map} beside it, and the two are read together: the hex is the program and
the map is the names. Delete one and the harness reports the other missing.

\subsection{Running the tests}
\label{c1:sec:tests}
From the repository root:

\begin{shell}
make test
\end{shell}

Every chapter's suite is cumulative, so this runs Chapter~1's now and will still be running it in
Chapter~5. What is in a suite decides for itself whether to run: the tests that check the pinned
device constants always run, and the ones that call your assembly are switched on by their
\code{.asm} file existing. Nothing has to be registered anywhere; the README beside the suite, in
\file{lectures/L01/exercises/test}, has the details.

Two failures worth recognising before you meet them:
\begin{itemize}
  \item \textbf{\code{Utils.SubroutinesAreDefined} fails and the file is plainly there.} The label
    is spelled differently from the name the test asks for, or the file did not assemble and the
    build said so further up. \code{avra} needs no \code{.global}, so it is never that.
  \item \textbf{A cycle count comes back as 100001.} That is the simulator's budget running out,
    which is what an infinite loop looks like from outside.
\end{itemize}

\subsection{Measuring a subroutine}
\label{c1:sec:measure}
A test tells you pass or fail. A cross-check needs the number itself, so the companion repository
ships a tool:

\begin{shell}
make measure SYMBOL=shift_bits ARG=5
\end{shell}

\begin{console}
shift_bits(r25:r24 = 0x0005, r22 = 0)
  cycles     40
  time       2.500 us at 16.0 MHz
  returned   r24 = 32 (0x20)

deepest stack: SP reached 0x08FD, 2 bytes below RAMEND
\end{console}

It assembles whatever is in \file{drivers/source}, loads it, sets \code{r25:r24} and \code{r22}
to the arguments you gave, calls the symbol, and reports what came back and what it cost.
\code{IMAGE=app} measures against the whole program instead of the library, which is what an
interrupt handler needs (\secref{c1:sec:images}).

\textbf{The figure excludes the \code{rcall} that would have got you there}, because there was no
\code{rcall}: the tool sets the program counter directly, the way the tests do. That is not a
defect and it is not hidden from you; it is the first term of the discrepancy that
Exercise~\ref{c1:ex:crosscheck}, the cross-check, asks you to account for.

\subsection{Opening this book's code in Microchip Studio}
\label{c1:sec:studio}
Everything in \file{drivers/source} is AVRASM2, which is what Microchip Studio's assembler takes.
A file from this book opens there and single-steps without being translated first, and that is
the reason for the choice: the Processor Status window showing \code{R16} change, and the I/O
view letting you tick a bit of \code{PINB} to fake a button press, apply to \textbf{your} code
rather than to a port of it.

Nothing in this book requires Studio, and no exercise depends on it. It is Windows-only, and
\code{make test} measures everything Studio would show you and rather more besides. But if you
have it, this is how it fits:
\begin{enumerate}
  \item \textbf{File $\rightarrow$ New $\rightarrow$ Project $\rightarrow$ AVR Assembler
    Project}, targeting the ATmega328P.
  \item Replace the generated \code{.asm} with one of yours, or add yours to the project.
  \item \textbf{Debug $\rightarrow$ Start Debugging and Break}, and pick \textbf{Simulator} as the
    tool when asked.
  \item \textbf{Debug $\rightarrow$ Windows $\rightarrow$ Processor Status} for the register file,
    and \textbf{I/O} for the peripherals.
\end{enumerate}

Two differences from the build here, neither of which is a syntax question:
\begin{itemize}
  \item \textbf{Studio assembles one file and this repository assembles all of them.}
    \file{ci/build.sh} generates a unit that includes every source (\secref{c1:sec:avra}); a Studio
    project has whatever you put in it. A subroutine that assembles here and is missing there is
    usually a file you did not add.
  \item \textbf{\code{m328Pdef.inc} comes from a different place.} Studio ships its own copy and
    finds it automatically; \code{avra} uses the one in \file{/usr/share/avra}. They agree about
    the ATmega328P, which is the only device this book targets.
\end{itemize}

\subsubsection{The vector table, in the units the datasheet uses}
\code{.org} counts \textbf{words} in AVRASM2, and so does the datasheet's vector table.
\code{.org 0x0006} is PCINT0 in both, and no conversion stands between them. Better still,
\code{m328Pdef.inc} names them:

\begin{asm}
.org PCI0addr
    rjmp isr_pcint0
\end{asm}

\code{PCI0addr}, \code{WDTaddr}, \code{OC1Aaddr} and the rest are defined in the device file, so
the magic number need not appear in your program at all. Chapter~3 is where this matters.

\begin{aside}
\textbf{Reading GNU \code{as} code elsewhere?} Most AVR assembly published online is written for
\code{avr-gcc}, where \code{.equ} takes a comma, register names come from \code{<avr/io.h>},
\code{PORTB} means the data space address \code{0x25} rather than the I/O address \code{0x05},
and \textbf{\code{.org} counts bytes}, so the same vector table is written \code{.org 0x000C}.
\secref{c6:sec:asmfromc} shows that dialect properly, because Chapter~6 is the one chapter where this book links
against a C compiler.
\end{aside}
